%% toc: P3 24 24.3
\chapter{Definable Sets with Parameters}
\label{sec:p03-w03-p03-003-definable-sets}

Inside one structure, a formula picks out the tuples that satisfy it.  When
some free variables of the formula are held fixed at chosen elements, the
\emph{parameters}, the remaining variables range over a set of tuples, and the
sets obtained in this way are the \emph{definable sets} of the structure.
They are the objects that most of model theory studies: Morley rank
(\cref{def:p03-morley-definable-rank}), strong minimality, algebraic closure
and the codes of imaginaries (\cref{def:p03-imaginaries-real-codes}) are all
properties of definable sets or are defined through them.  This chapter gives
the general account for tuples of any finite length and parameters from a
prescribed set.

\emph{Inputs.}  A first-order signature \(\Lang\), an \(\Lang\)-structure
\(\mathcal M\) with domain \(M\), a set \(A\subseteq M\) of admissible
parameters (possibly empty), and formulas of \(\Lang\) together with a
declared list of their free variables.

\emph{Conclusions.}  For each \(n\geq1\) the \(A\)-definable subsets of
\(M^n\) form a Boolean algebra of sets; the whole family is closed under
taking images and preimages along maps that permute, delete, repeat or add
coordinates, and under fixing coordinates at elements of \(A\); and it is the
least family with these closure properties containing the basic relations of
\(\mathcal M\) (\cref{thm:p03-w03-p03-003-closure,thm:p03-w03-p03-003-least-family}).
Every automorphism of \(\mathcal M\) fixing \(A\) pointwise maps each
\(A\)-definable set onto itself, although invariance does not conversely
imply definability (\cref{thm:p03-w03-p03-003-automorphism,thm:p03-w03-p03-003-rigid}).
Passing to an elementary extension \(\mathcal N\) extends each \(A\)-definable
set of \(\mathcal M\) to an \(A\)-definable set of \(\mathcal N\) with the same
trace on \(M\), and this extension is an isomorphism of the two Boolean
algebras (\cref{thm:p03-w03-p03-003-elementary-extension}); for a substructure
that is not elementary the transfer fails
(\cref{def:p03-w03-p03-003-substructure-failure}).

The definition of definable and \(A\)-definable sets follows
\cite[p.~17]{Marker2000IntroductionModelTheory}, and the characterization by
closure conditions extends \cite[Proposition~1.3, p.~18]{Marker2000IntroductionModelTheory},
which is stated there without proof and for arbitrary parameters only.  All
proofs are written in Annex~\ref{sec:annex-p03-w03-p03-003-definable-sets}.
Satisfaction, assignments and free variables are those of
\cref{sec:fol-semantics-02}; the variable set is the countable set fixed in
\cref{sec:fol-syntax-semantics-01}, and a constant symbol is a function symbol
of arity zero.

\section{Formulas with declared variables and coordinate maps}
\label{sec:p03-w03-p03-003-declared}

\begin{hbconvention}{def:p03-w03-p03-003-declared-formula}{Declared formulas and their solution sets}
Let \(\mathcal M\) be an \(\Lang\)-structure with domain \(M\).
\begin{enumerate}
 \item For a list \(\bar v=(v_1,\ldots,v_k)\) of pairwise distinct variables,
       an assignment \(s\) and \(\bar c=(c_1,\ldots,c_k)\in M^k\), write
       \(s[\bar v\mapsto\bar c]\) for the assignment that sends \(v_i\) to
       \(c_i\) and agrees with \(s\) elsewhere.  Because the \(v_i\) are
       distinct, this equals the result of the successive single updates
       \(s[v_1\mapsto c_1]\cdots[v_k\mapsto c_k]\) taken in any order.
 \item A \emph{declared formula} \(\varphi(x_1,\ldots,x_n;y_1,\ldots,y_m)\),
       abbreviated \(\varphi(\bar x;\bar y)\), is an \(\Lang\)-formula
       \(\varphi\) together with a list \((\bar x;\bar y)\) of \(n+m\)
       pairwise distinct variables such that
       \(\FV(\varphi)\subseteq\{x_1,\ldots,x_n,y_1,\ldots,y_m\}\).  Here
       \(n\geq1\) and \(m\geq0\); the \(x_i\) are the \emph{solution
       variables} and the \(y_j\) the \emph{parameter variables}.  Not every
       listed variable need occur free in \(\varphi\).
 \item For \(\bar a\in M^n\) and \(\bar b\in M^m\), write
       \(\mathcal M\entails\varphi(\bar a,\bar b)\) when
       \(\mathcal M,s[\bar x\mapsto\bar a,\bar y\mapsto\bar b]\entails\varphi\)
       for some assignment \(s\).  By \cref{lem:fol-coincidence} the choice
       of \(s\) is immaterial, so ``some'' may be replaced by ``every''.
 \item The \emph{solution set} of \(\varphi(\bar x;\bar y)\) at the parameter
       tuple \(\bar b\in M^m\) is
       \[
        \varphi(M^n,\bar b)
        \defeq\{\bar a\in M^n:\mathcal M\entails\varphi(\bar a,\bar b)\}.
       \]
       When \(m=0\) it is written \(\varphi(M^n)\).
\end{enumerate}
The solution set depends on the declared list and not only on the formula:
with the binary relation symbol \(<\), the formula \(x<y\) declared as
\((x;y)\) at the parameter \(b\) gives the elements below \(b\), while the
same formula declared as \((y;x)\) gives the elements above \(b\).  The list
also fixes the order of coordinates in \(M^n\).
\end{hbconvention}

\begin{hbdefinition}{def:p03-w03-p03-003-coordinate-maps}{Coordinate maps}
For integers \(k,n\geq1\) and a function
\(\tau:\{1,\ldots,k\}\to\{1,\ldots,n\}\), the \emph{coordinate map} of
\(\tau\) on \(M\) is
\[
 \tau_M:M^n\to M^k,\qquad
 \tau_M(a_1,\ldots,a_n)=(a_{\tau(1)},\ldots,a_{\tau(k)}).
\]
For \(X\subseteq M^n\) its \emph{image} is
\(\tau_M(X)=\{\tau_M(\bar a):\bar a\in X\}\), and for \(Y\subseteq M^k\) its
\emph{preimage} is
\(\tau_M^{-1}(Y)=\{\bar a\in M^n:\tau_M(\bar a)\in Y\}\).  The named
operations on sets of tuples are special cases:
\begin{enumerate}
 \item \emph{permutation of coordinates}: \(k=n\) and \(\tau\) a bijection;
 \item \emph{projection} onto the first \(k\) coordinates: \(k\leq n\) and
       \(\tau(i)=i\), so that
       \(\tau_M(X)=\{(a_1,\ldots,a_k):(a_1,\ldots,a_n)\in X
       \text{ for some }a_{k+1},\ldots,a_n\in M\}\);
 \item \emph{identification of coordinates}: \(\tau\) not injective; for
       instance \(k=2\), \(n=1\), \(\tau(1)=\tau(2)=1\) gives
       \(\tau_M^{-1}(Y)=\{a\in M:(a,a)\in Y\}\);
 \item \emph{adding dummy coordinates}: \(k<n\) and \(\tau(i)=i\), so that
       \(\tau_M^{-1}(Y)=Y\times M^{n-k}\).
\end{enumerate}
Projection onto any other set of coordinates is a permutation followed by
projection onto the first coordinates.
\end{hbdefinition}

\section{\texorpdfstring{\(A\)}{A}-definable sets and relations}
\label{sec:p03-w03-p03-003-a-definable-section}

\begin{hbdefinition}{def:p03-w03-p03-003-a-definable}{Definable sets and relations over a parameter set}
Let \(\mathcal M\) be an \(\Lang\)-structure, \(A\subseteq M\) and
\(n\geq1\).  A set \(X\subseteq M^n\) is \emph{\(A\)-definable}, or
\emph{definable over \(A\)}, in \(\mathcal M\) iff there are a declared
formula \(\varphi(x_1,\ldots,x_n;y_1,\ldots,y_m)\) and a tuple
\(\bar b\in A^m\) with
\[
 X=\varphi(M^n,\bar b).
\]
The formula, with its declared list, and the tuple \(\bar b\) together
\emph{define} \(X\) over \(A\).  Write
\(\operatorname{Def}_n(\mathcal M/A)\) for the set of \(A\)-definable subsets
of \(M^n\).  A set is \emph{\(\varnothing\)-definable}, or definable without
parameters, iff it is \(A\)-definable for \(A=\varnothing\), which forces
\(m=0\); it is \emph{definable} iff it is \(M\)-definable.  An
\emph{\(A\)-definable \(n\)-ary relation} on \(M\) is an \(A\)-definable
subset of \(M^n\); the two terms name the same objects.

These are the notions of \cite[p.~17]{Marker2000IntroductionModelTheory},
where the formula is written \(\phi(x_1,\ldots,x_{n+m})\) with the parameter
variables last and the case \(m=0\) is called \(\varnothing\)-definable.
\end{hbdefinition}

Every \(A\)-definable set is \(A_0\)-definable for the finite set \(A_0\) of
entries of a defining tuple \(\bar b\), and \(A\subseteq B\subseteq M\)
gives \(\operatorname{Def}_n(\mathcal M/A)\subseteq
\operatorname{Def}_n(\mathcal M/B)\).  Thus each definable set uses only
finitely many parameters even when \(A\) is infinite.

\begin{hblemma}{thm:p03-w03-p03-003-normal-form}{proved=annex:p03-w03-p03-003-normal-form}{Renaming, padding and constants for parameters}
Let \(\varphi(x_1,\ldots,x_n;y_1,\ldots,y_m)\) be a declared
\(\Lang\)-formula.
\begin{enumerate}
 \item \emph{Renaming.}  For every list \((\bar u;\bar w)\) of \(n+m\)
       pairwise distinct variables there is an \(\Lang\)-formula \(\varphi'\)
       with \(\FV(\varphi')\subseteq\{\bar u,\bar w\}\) such that, for every
       \(\Lang\)-structure \(\mathcal K\) and every \(\bar b\in K^m\), the
       declared formula \(\varphi'(\bar u;\bar w)\) has
       \(\varphi'(K^n,\bar b)=\varphi(K^n,\bar b)\).  The formula
       \(\varphi'\) is obtained from \(\varphi\) and the two lists alone.
 \item \emph{Padding.}  If \(\bar z\) is a list of \(r\) pairwise distinct
       variables outside \(\operatorname{Var}(\varphi)\cup
       \{\bar x,\bar y\}\), then \(\varphi\) declared as
       \((\bar x;\bar y\bar z)\) satisfies
       \(\varphi(K^n,\bar b\bar c)=\varphi(K^n,\bar b)\) for every
       \(\Lang\)-structure \(\mathcal K\), \(\bar b\in K^m\) and
       \(\bar c\in K^r\).
 \item \emph{Constants for parameters.}  Let \(A\subseteq M\) and let
       \(\mathcal M_A\) be the \(\Lang(A)\)-expansion of
       \cref{def:parameter-expansion}.  A set \(X\subseteq M^n\) is
       \(A\)-definable in \(\mathcal M\) iff \(X=\psi(M^n)\) for a declared
       \(\Lang(A)\)-formula \(\psi(x_1,\ldots,x_n;)\) evaluated in
       \(\mathcal M_A\), that is, iff \(X\) is \(\varnothing\)-definable in
       \(\mathcal M_A\).
\end{enumerate}
\end{hblemma}

By the renaming clause, any finite family of definable sets in the same
\(M^n\) can be given defining formulas with one common list of solution
variables and pairwise disjoint lists of parameter variables.  The proofs of
the closure results use this repeatedly.

\begin{hbtheorem}{thm:p03-w03-p03-003-closure}{proved=annex:p03-w03-p03-003-closure}{Closure of the \(A\)-definable family}
Let \(\mathcal M\) be an \(\Lang\)-structure, \(A\subseteq M\) and
\(k,n,r\geq1\).
\begin{enumerate}
 \item \emph{Basic sets.}  The following are \(\varnothing\)-definable, hence
       \(A\)-definable: \(M^n\); the relation \(R^{\mathcal M}\subseteq M^p\)
       for each relation symbol \(R\) of arity \(p\geq1\); the graph
       \(\{(\bar a,f^{\mathcal M}(\bar a)):\bar a\in M^p\}\subseteq M^{p+1}\)
       of each function symbol \(f\) of arity \(p\geq0\), which for a
       constant \(c\) is \(\{c^{\mathcal M}\}\subseteq M^1\); and each
       diagonal \(\Delta^n_{ij}=\{\bar a\in M^n:a_i=a_j\}\) for
       \(1\leq i,j\leq n\).
 \item \emph{Boolean operations.}  \(\operatorname{Def}_n(\mathcal M/A)\)
       contains \(\varnothing\) and \(M^n\) and is closed under complement
       in \(M^n\) and under finite intersections and unions; it is a Boolean
       subalgebra of the power set of \(M^n\).
 \item \emph{Preimages.}  If \(Y\in\operatorname{Def}_k(\mathcal M/A)\) and
       \(\tau:\{1,\ldots,k\}\to\{1,\ldots,n\}\), then
       \(\tau_M^{-1}(Y)\in\operatorname{Def}_n(\mathcal M/A)\).  In
       particular the family is closed under permutation and identification
       of coordinates and under adding dummy coordinates.
 \item \emph{Images.}  If \(X\in\operatorname{Def}_n(\mathcal M/A)\) and
       \(\tau:\{1,\ldots,k\}\to\{1,\ldots,n\}\), then
       \(\tau_M(X)\in\operatorname{Def}_k(\mathcal M/A)\).  In particular the
       family is closed under projection.
 \item \emph{Fibres over parameters.}  If
       \(X\in\operatorname{Def}_{n+r}(\mathcal M/A)\) and
       \(\bar c\in A^r\), then
       \(\{\bar a\in M^n:(\bar a,\bar c)\in X\}
       \in\operatorname{Def}_n(\mathcal M/A)\).
 \item \emph{Uniformity.}  In each of (2)--(5) the annex constructs, from
       defining formulas and parameter tuples of the given sets, a declared
       formula and a parameter tuple that define the resulting set.  The
       construction depends only on the given formulas, lists, tuples and
       \(\tau\), and the same formula and tuple define the corresponding
       resulting set in every \(\Lang\)-structure in which the given
       formulas and tuples are evaluated.
\end{enumerate}
\end{hbtheorem}

\begin{hbtheorem}{thm:p03-w03-p03-003-least-family}{proved=annex:p03-w03-p03-003-least}{The \(A\)-definable sets form the least closed family}
Let \(\mathcal M\) be an \(\Lang\)-structure and \(A\subseteq M\).  Let
\(\mathcal D=(\mathcal D_n)_{n\geq1}\), with
\(\mathcal D_n\subseteq\mathcal P(M^n)\), be the least family such that
\begin{enumerate}
 \item every basic set listed in
       \cref{thm:p03-w03-p03-003-closure}(1) belongs to the corresponding
       \(\mathcal D_n\);
 \item each \(\mathcal D_n\) is closed under complement in \(M^n\), finite
       intersection and finite union;
 \item \(\tau_M^{-1}(Y)\in\mathcal D_n\) whenever \(Y\in\mathcal D_k\) and
       \(\tau:\{1,\ldots,k\}\to\{1,\ldots,n\}\);
 \item \(\tau_M(X)\in\mathcal D_k\) whenever \(X\in\mathcal D_n\) and
       \(\tau:\{1,\ldots,k\}\to\{1,\ldots,n\}\); and
 \item \(\{\bar a\in M^n:(\bar a,\bar c)\in X\}\in\mathcal D_n\) whenever
       \(X\in\mathcal D_{n+r}\) and \(\bar c\in A^r\).
\end{enumerate}
Then \(\mathcal D_n=\operatorname{Def}_n(\mathcal M/A)\) for every
\(n\geq1\).  The least family exists: it is the componentwise intersection
of all families with properties (1)--(5), and the family of all subsets of
the powers \(M^n\) is one such family.
\end{hbtheorem}

With \(A=M\) this is \cite[Proposition~1.3, p.~18]{Marker2000IntroductionModelTheory}
up to one difference of bookkeeping.  Marker's list has a clause for the
graphs of the functions of the structure and none for its distinguished
elements.  When \(A=M\) nothing is lost: the singleton
\(\{c^{\mathcal M}\}\) is the fibre of the diagonal \(\Delta^2_{12}\) over
the parameter \(c^{\mathcal M}\), so clause (5) already supplies it.  For a
smaller \(A\) the constant clause in (1) is needed, since \(c^{\mathcal M}\)
need not belong to \(A\).

\begin{hbmethod}{def:p03-w03-p03-003-procedure}{Writing a defining formula for a constructed set}
Given defining data for the input sets, a defining formula for a set built by
\cref{thm:p03-w03-p03-003-closure} is written as follows.  Each step lists
its side condition.
\begin{enumerate}
 \item \emph{Separate the variables.}  Rename every input formula by
       \cref{thm:p03-w03-p03-003-normal-form}(1) so that all inputs living in
       the same \(M^n\) share one list \(\bar x\) of solution variables and
       their parameter lists are pairwise disjoint and disjoint from
       \(\bar x\).
 \item \emph{Boolean step.}  Join the renamed formulas by \(\neg\),
       \(\land\), \(\lor\); concatenate the parameter lists and tuples.
       Side condition: the parameter lists must be disjoint, otherwise one
       variable would be asked to carry two parameter values.
 \item \emph{Preimage under \(\tau\).}  From
       \(\varphi(x_1,\ldots,x_k;\bar y)\) form
       \(\exists x_1\cdots\exists x_k\bigl(\bigwedge_{i\leq k}x_i=u_{\tau(i)}
       \land\varphi\bigr)\), declared as \((u_1,\ldots,u_n;\bar y)\).
       Side condition: the \(u_l\) are distinct and lie outside
       \(\operatorname{Var}(\varphi)\cup\{\bar x,\bar y\}\).
 \item \emph{Image under \(\tau\).}  From \(\varphi(x_1,\ldots,x_n;\bar y)\)
       form
       \(\exists x_1\cdots\exists x_n\bigl(\bigwedge_{i\leq k}u_i=x_{\tau(i)}
       \land\varphi\bigr)\), declared as \((u_1,\ldots,u_k;\bar y)\), with
       the same side condition on the \(u_i\).  For projection onto the first
       \(k<n\) coordinates the simpler formula
       \(\exists x_{k+1}\cdots\exists x_n\,\varphi\), declared as
       \((x_1,\ldots,x_k;\bar y)\), also works.
 \item \emph{Fibre.}  To fix the last \(r\) solution variables at
       \(\bar c\in A^r\), move them to the front of the parameter list and
       put \(\bar c\) in front of the parameter tuple; the formula is
       unchanged.
\end{enumerate}
No step changes the parameter tuple except by concatenation and by
inserting elements of \(A\), so the result is again defined over \(A\).
\end{hbmethod}

\section{Automorphisms and parameters}
\label{sec:p03-w03-p03-003-automorphisms}

For \(A\subseteq M\), write \(\operatorname{Aut}(\mathcal M/A)\) for the set
of automorphisms \(\sigma\) of \(\mathcal M\) with \(\sigma(a)=a\) for every
\(a\in A\).  An automorphism acts on tuples coordinatewise and on sets of
tuples by images: \(\sigma(X)=\{\sigma\bar a:\bar a\in X\}\).

\begin{hbtheorem}{thm:p03-w03-p03-003-automorphism}{proved=annex:p03-w03-p03-003-auto}{Isomorphism transport and automorphism invariance}
Let \(h:\mathcal M\to\mathcal N\) be an isomorphism of \(\Lang\)-structures,
\(\varphi(\bar x;\bar y)\) a declared formula with \(n\) solution and \(m\)
parameter variables, and \(\bar b\in M^m\).  Then
\[
 h\bigl(\varphi(M^n,\bar b)\bigr)=\varphi(N^n,h\bar b).
\]
Consequently, if \(A\subseteq M\), \(X\in\operatorname{Def}_n(\mathcal M/A)\)
and \(\sigma\in\operatorname{Aut}(\mathcal M/A)\), then \(\sigma(X)=X\).  In
contrapositive form: if some \(\sigma\in\operatorname{Aut}(\mathcal M/A)\)
has \(\sigma(X)\neq X\), then \(X\) is not \(A\)-definable.
\end{hbtheorem}

The theorem rests on isomorphism invariance of satisfaction,
\cref{lem:iso-invariance}, which is \cite[Proposition~1.2, p.~17]{Marker2000IntroductionModelTheory}
in Marker's account.  Invariance is a necessary condition for
\(A\)-definability only.

\begin{hbtheorem}{thm:p03-w03-p03-003-rigid}{proved=annex:p03-w03-p03-003-rigid}{An invariant set that is not definable}
Let \(\Lang=\{<\}\) with \(<\) a binary relation symbol, and let
\(\mathcal M=(\N,<)\) be the natural numbers with their usual order.
\begin{enumerate}
 \item The only automorphism of \(\mathcal M\) is the identity.
 \item Only countably many subsets of \(\N\) are definable in \(\mathcal M\),
       with arbitrary parameters.
 \item Hence some \(X\subseteq\N\) is not definable in \(\mathcal M\) even
       though \(\sigma(X)=X\) for every \(\sigma\in\operatorname{Aut}
       (\mathcal M/A)\) and every \(A\subseteq\N\).
\end{enumerate}
\end{hbtheorem}

A converse of the invariance theorem therefore needs extra hypotheses on the
structure.  Marker gives one at the level of single tuples rather than sets:
two tuples with the same complete type over \(A\) are mapped to one another
by an automorphism fixing \(A\) of some elementary extension
\cite[Proposition~5.1, p.~28]{Marker2000IntroductionModelTheory}, and inside
\(\mathcal M\) itself when \(\mathcal M\) is saturated and \(\lvert A\rvert<
\lvert M\rvert\) \cite[Lemma~6.2, p.~30]{Marker2000IntroductionModelTheory}.

\section{Definable sets under elementary extension}
\label{sec:p03-w03-p03-003-extension}

Elementary substructures and elementary extensions are those of
\cref{sec:maps-elementarity-03}: \(\mathcal M\preccurlyeq\mathcal N\) means
that \(\mathcal M\subseteq\mathcal N\) and every formula has the same truth
value in both structures at every tuple from \(M\).

\begin{hbtheorem}{thm:p03-w03-p03-003-elementary-extension}{proved=annex:p03-w03-p03-003-elem}{Transfer of definable sets to an elementary extension}
Let \(\mathcal M\preccurlyeq\mathcal N\), let \(A\subseteq M\) and
\(n\geq1\).
\begin{enumerate}
 \item \emph{Trace.}  For every declared formula \(\varphi(\bar x;\bar y)\)
       and \(\bar b\in M^m\),
       \(\varphi(N^n,\bar b)\cap M^n=\varphi(M^n,\bar b)\).
 \item \emph{Independence of the definition.}  If declared formulas
       \(\varphi(\bar x;\bar y)\) and \(\psi(\bar x';\bar y')\) with tuples
       \(\bar b\) and \(\bar c\) from \(M\) define the same subset of
       \(M^n\), then \(\varphi(N^n,\bar b)=\psi(N^n,\bar c)\).
       Hence for a definable \(X\subseteq M^n\) the set
       \(X^{\mathcal N}\defeq\varphi(N^n,\bar b)\), for any defining data
       of \(X\) with parameters in \(M\), is well defined.
 \item \emph{Isomorphism of the algebras.}  The map
       \(X\mapsto X^{\mathcal N}\) sends
       \(\operatorname{Def}_n(\mathcal M/A)\) bijectively onto
       \(\operatorname{Def}_n(\mathcal N/A)\), with
       \(X^{\mathcal N}\cap M^n=X\).  It sends \(M^n\) to \(N^n\) and
       commutes with complement, finite intersection and finite union; it
       preserves and reflects inclusion; and it commutes with the preimages,
       images and fibres of \cref{thm:p03-w03-p03-003-closure}:
       \((\tau_M^{-1}(Y))^{\mathcal N}=\tau_N^{-1}(Y^{\mathcal N})\),
       \((\tau_M(X))^{\mathcal N}=\tau_N(X^{\mathcal N})\), and the fibre of
       \(X^{\mathcal N}\) over \(\bar c\in A^r\) is the extension of the
       fibre of \(X\) over \(\bar c\).
 \item \emph{Finite and empty sets.}  For a definable
       \(X\subseteq M^n\), \(X=\varnothing\) iff
       \(X^{\mathcal N}=\varnothing\), and \(X\) is finite iff
       \(X^{\mathcal N}\) is finite, in which case \(X^{\mathcal N}=X\).
\end{enumerate}
\end{hbtheorem}

An infinite definable set may acquire new elements in an elementary
extension: \(X^{\mathcal N}\) contains \(X\) and agrees with it on \(M^n\),
but may contain tuples with entries outside \(M\).  Clause (3) says that the
\(A\)-definable sets of \(\mathcal N\) are exactly these extensions, so
passing to an elementary extension adds no new \(A\)-definable sets; sets
definable in \(\mathcal N\) with parameters outside \(M\) are not covered.
The parameter-name expansion of \cref{thm:p03-w03-p03-003-normal-form}(3)
connects the theorem with the elementary diagram of \cref{def:diagrams}:
\(\mathcal M\preccurlyeq\mathcal N\) says that \(\mathcal N_M\) satisfies
every \(\Lang(M)\)-sentence true in \(\mathcal M_M\).

\begin{hbwarning}{def:p03-w03-p03-003-substructure-failure}{Transfer fails for a substructure that is not elementary}
Clauses (1) and (2) of \cref{thm:p03-w03-p03-003-elementary-extension} need
elementarity, not only \(\mathcal M\subseteq\mathcal N\).  Take
\(\Lang=\{<\}\), \(\mathcal M=(\N,<)\) and \(\mathcal N=(\mathbb Z,<)\), so
that \(\mathcal M\subseteq\mathcal N\) as in the example of
\cref{sec:maps-elementarity-02}.  Declare
\(\varphi(x;)\defeq\exists y\,(y<x)\) and
\(\psi(x;y)\defeq\neg\,x=y\).  In \(\mathcal M\), \(\varphi(\N)\) is the set
of positive natural numbers, because \(0\) is the only element with nothing
below it; and \(\psi(\N,0)\) is the same set.  In \(\mathcal N\), every
integer has a smaller one, so \(\varphi(\mathbb Z)=\mathbb Z\), while
\(\psi(\mathbb Z,0)=\mathbb Z\setminus\{0\}\).  Thus
\(\varphi(\mathbb Z)\cap\N=\N\neq\varphi(\N)\), so the trace clause fails,
and two definitions of one subset of \(\N\) have different solution sets in
\(\mathbb Z\), so no extension map \(X\mapsto X^{\mathcal N}\) is
determined.
\end{hbwarning}

\section{Worked example: the rational order}
\label{sec:p03-w03-p03-003-example}

\begin{hbexample}{def:p03-w03-p03-003-rational-order}{Definable subsets of the rational line}
Let \(\Lang=\{<\}\) and \(\mathcal Q=(\mathbb Q,<)\).  For rationals \(q\)
and \(r>0\), the maps \(x\mapsto x+q\) and \(x\mapsto rx\) are bijections of
\(\mathbb Q\) onto itself with inverses \(x\mapsto x-q\) and
\(x\mapsto r^{-1}x\), and \(a<b\) iff \(a+q<b+q\) iff \(ra<rb\); so both are
automorphisms of \(\mathcal Q\).

\emph{No parameters.}  Let \(X\in\operatorname{Def}_1(\mathcal Q/\varnothing)\)
be nonempty, say \(a\in X\), and let \(b\in\mathbb Q\).  The translation
\(x\mapsto x+(b-a)\) lies in \(\operatorname{Aut}(\mathcal Q/\varnothing)\)
and sends \(a\) to \(b\), so \(b\in X\) by
\cref{thm:p03-w03-p03-003-automorphism}.  Hence
\(\operatorname{Def}_1(\mathcal Q/\varnothing)=\{\varnothing,\mathbb Q\}\),
both being definable, by \(\neg\,x=x\) and \(x=x\).

\emph{One parameter.}  Let \(A=\{0\}\).  Declare \(\theta(x;y)\defeq y<x\),
\(\theta'(x;y)\defeq x<y\) and \(\theta''(x;y)\defeq x=y\).  At the parameter
\(0\) they define \(P=\{a:a>0\}\), \(N=\{a:a<0\}\) and \(\{0\}\).  If
\(a,b\) are both positive or both negative, then \(r=b/a>0\) and
\(x\mapsto rx\) lies in \(\operatorname{Aut}(\mathcal Q/\{0\})\) and sends
\(a\) to \(b\).  So every \(\{0\}\)-definable \(X\subseteq\mathbb Q\) that
meets \(P\) contains \(P\), and likewise for \(N\); hence \(X\) is a union of
some of \(P\), \(N\), \(\{0\}\).  Conversely each of the eight such unions is
\(\{0\}\)-definable by the Boolean clause of
\cref{thm:p03-w03-p03-003-closure}.  Thus
\(\operatorname{Def}_1(\mathcal Q/\{0\})\) has exactly eight members, and
\(P\) is \(\{0\}\)-definable but not \(\varnothing\)-definable.

\emph{Coordinate operations.}  Let
\(L=\{(a_1,a_2):a_1<a_2\}\in\operatorname{Def}_2(\mathcal Q/\varnothing)\).
Its projection onto the first coordinate is \(\mathbb Q\), since no rational
is largest; its preimage under the identification
\(\tau(1)=\tau(2)=1\) is \(\{a:a<a\}=\varnothing\); its image under the
transposition is \(\{(a_1,a_2):a_2<a_1\}\); and its fibre over the parameter
\(0\) in the second coordinate is \(N\), defined by \(x_1<x_2\) declared as
\((x_1;x_2)\) at \(0\), exactly as in step (5) of
\cref{def:p03-w03-p03-003-procedure}.
\end{hbexample}

\section{Relationships and limits}
\label{sec:p03-w03-p03-003-relationships}

\emph{Definable sets and elementary classes.}  An elementary class
(\cref{def:p03-elementary-class}) is a class of whole structures cut out by
sentences; a definable set is a set of tuples inside one fixed structure cut
out by one formula with parameters.  The two notions meet in
\cref{thm:p03-w03-p03-003-normal-form}(3): an \(A\)-definable set of
\(\mathcal M\) is a \(\varnothing\)-definable set of the expansion
\(\mathcal M_A\), which is how the diagrams of \cref{def:diagrams} turn
statements with parameters into statements about classes of expanded
structures.

\emph{Definable sets and explicit definitions.}  Explicit definability in
the sense of \cref{def:p03-implicit-explicit-definability} asks for a single
parameter-free formula that defines a new relation symbol in every model of
a theory.  In a model of the theory with domain \(M\), an explicitly
defined symbol \(P\) of arity \(n\geq1\) is
interpreted by the \(\varnothing\)-definable set \(\theta(M^n)\) of its
\(\Lang\)-reduct, where \(\theta\) is the defining formula declared with
solution variables \(\bar x\); one formula serves every model of the theory.

\emph{Interpretations, ranks and codes.}  The domain formula and relation
formulas of a relative interpretation
(\cref{def:p03-vaught-relative-interpretation}) are declared formulas
without parameters, so the interpreted structure lives on
\(\varnothing\)-definable sets of the target model.  Morley rank
(\cref{def:p03-morley-definable-rank}) and strong minimality are defined on
definable sets with arbitrary parameters, and the elementary extensions used
there act on those sets as described in
\cref{thm:p03-w03-p03-003-elementary-extension}.  The definable closure of
\cref{def:p03-imaginaries-real-codes} is built from formulas over finite
subsets of a parameter set, in agreement with the finite-character remark
following \cref{def:p03-w03-p03-003-a-definable}.

\emph{Limits.}  Definability depends on the signature, on the structure and
on \(A\): \cref{def:p03-w03-p03-003-rational-order} exhibits a set that is
definable over \(\{0\}\) and not over \(\varnothing\).  Automorphism
invariance does not characterize definability
(\cref{thm:p03-w03-p03-003-rigid}), and the transfer to larger structures
holds for elementary extensions and fails for arbitrary extensions
(\cref{def:p03-w03-p03-003-substructure-failure}).  The account here is for
one-sorted first-order logic with finitely many free variables per formula;
sets of infinite tuples and definability in other logics are not covered by
these theorems.

\emph{Source.}  The definitions of definable, \(A\)-definable and
\(\varnothing\)-definable sets are those of
\cite[p.~17]{Marker2000IntroductionModelTheory}, read in the open-access
copy of the published chapter; the closure characterization for \(A=M\) is
\cite[Proposition~1.3, p.~18]{Marker2000IntroductionModelTheory}.  The
\(A\)-relative closure theorem, the normal-form lemma, the automorphism and
elementary-extension theorems and the two examples are given here with
proofs in Annex~\ref{sec:annex-p03-w03-p03-003-definable-sets}.
